% !PRINTED 524
\weberpage{524}
%<<CONT
und
\weq{7}{\begin{gathered}
v = \alpha + \varepsilon\alpha_1 + \varepsilon^2\alpha_2 \\
v' = \alpha + \varepsilon^2\alpha_1 + \varepsilon\alpha_2;
\end{gathered}}
dann ist der Erfolg von $\pi$ der, dass $v$ in $\varepsilon^2 v$, $v'$ in $\varepsilon v'$ übergeht, während $v^3$, $v'^3$ und $vv'$ ungeändert bleiben, also in dem Körper $\Omega'$ enthalten sind. Es hat keine Schwierigkeit, diese Grössen zu berechnen. Setzen wir zur Abkürzung
\[
\begin{gathered}
A = \alpha^2\alpha_1 + \alpha_1^2\alpha_2 + \alpha_2^2\alpha, \\
A' = \alpha_1^2\alpha + \alpha_2^2\alpha_1 + \alpha^2\alpha_2,
\end{gathered}
\]
so ist
\[
\begin{gathered}
A + A' = ab - 3c \\
A - A' = \sqrt{D},
\end{gathered}
\]
und für $v^3$ erhält man
\[
v^3 = \alpha^3 + \alpha_1^3 + \alpha_2^3 + 6\alpha\alpha_1\alpha_2 + 3\varepsilon A + 3\varepsilon^2 A',
\]
also
\weq{8}{v^3 = a^3 - \frac{9}{2}ab + \frac{27}{2}c + \frac{3}{2}\sqrt{-3D},}
und ebenso
\weq{9}{v'^3 = a^3 - \frac{9}{2}ab + \frac{27}{2}c - \frac{3}{2}\sqrt{-3D}.}

Es ist aber auch
\weq{10}{vv' = \alpha^2 + \alpha_1^2 + \alpha_2^2 - \alpha\alpha_1 - \alpha\alpha_2 - \alpha_1\alpha_2 = a^2 - 3b,}
und aus (8), (9), (10) erhält man in Uebereinstimmung mit einer früheren Formel [§. 47, (8)]
\[
27D = 4(a^2 - 3b)^3 - (2a^3 - 9ab + 27c)^2.
\]
Fügt man zu (7) noch die Gleichung
\[
a = \alpha + \alpha_1 + \alpha_2,
\]
so findet man in Uebereinstimmung mit der \gesperrt{Cardanischen} Formel
\[
\begin{gathered}
3\alpha = a + v + v' \\
3\alpha_1 = a + \varepsilon^2 v + \varepsilon v' \\
3\alpha_2 = a + \varepsilon v + \varepsilon^2 v'.
\end{gathered}
\]

\begin{center}
§. 160.

\gesperrt{Permutationsgruppen von vier Elementen.}
\end{center}

Ein gutes Beispiel für die \gesperrt{Galois'sche} Theorie bieten die Gleichungen vierten Grades, wo die Verhältnisse so einfach
%CONT>>
